\section{Introduction}
\label{sec:introduction}

\IEEEPARstart{M}{ultidimensional} signals often contain strong structure across
space, spectrum, time, view, or acquisition mode. Low-rank tensor models
exploit this structure by representing a signal through a small number of
coupled low-dimensional factors. They have become standard tools for signal
representation, denoising, completion, and recovery
\cite{Kolda2009Tensor,Sidiropoulos2017TSP,Liu2013TensorCompletion,
LiuMoitra2020TensorCompletion}. Tucker decomposition is especially useful in
this setting because it separates a compact core from mode-wise factors and
provides an explicit multilinear rank budget.

In many sensing problems, however, the structured quantity of interest is not
recorded directly. Camera response curves, radiometric processing, saturation,
reflectance effects, and post-nonlinear hyperspectral mixing are standard
examples in which a physical signal is transformed before it is stored as data
\cite{Debevec1997HDR,Grossberg2004CameraResponse,
BioucasDias2012HyperspectralUnmixing,Altmann2012PostNonlinear,
Heylen2014NonlinearUnmixing}. Related ideas also appear in generalized
low-rank and single-index matrix models, where low-dimensional structure is
observed through a non-Gaussian or nonlinear entrywise map
\cite{Udell2016GLRM,Ganti2015SingleIndex}. These observation models share a
simple lesson: a compact pre-measurement signal can look more complicated after
the measurement process.

For Tucker recovery, this distinction matters because an entrywise nonlinear
response can inflate the apparent multilinear rank of the measured tensor. An
underlying low-rank signal can therefore appear to require a substantially
larger Tucker rank after measurement. From a signal recovery perspective, this
creates a model mismatch: the rank budget selected for the underlying signal
may be appropriate, while the linear Tucker observation model is not.

One way to reduce this mismatch is to increase the Tucker rank. This is often
effective, but it changes the entire multilinear representation: the core
grows rapidly, each mode factor gains additional degrees of freedom, and the
interpretation of a fixed rank budget for the underlying signal is lost. A second approach is to
replace the multilinear model with a more flexible nonlinear tensor model, such
as a kernel, Gaussian-process, generalized-link, or neural tensor decoder
\cite{Luo2024LRTFR,saragadam2024deeptensor,ahn2024neural}. Such methods are
powerful, but their flexibility can obscure the relation between the recovered
signal and a chosen Tucker representation, and may introduce additional training and
model-selection burdens.

This paper studies a narrower and more structured alternative. We assume that
the underlying signal is well represented by a Tucker tensor, but that the
measurements are generated from that underlying signal through an unknown
entrywise response. We approximate this response in a finite dictionary of
scalar functions and learn the response jointly with the Tucker core and
factors. If
\begin{equation}
  \cS =
  \tucker{\cX}{\bA^{(1)},\ldots,\bA^{(N)}}
\end{equation}
is the underlying Tucker signal, the measured tensor is modeled as
\begin{equation}
  \widehat{\cY}
  =
  f_{\bbeta}(\cS),
  \label{eq:intro_ntdpl}
\end{equation}
where $f_{\bbeta}(s)=\sum_{q=0}^{Q}\beta_q\psi_q(s)$ for a chosen scalar
dictionary $\{\psi_q\}_{q=0}^{Q}$. We call this a dictionary-linked Tucker
model. It keeps the Tucker ranks fixed and adds only the response coefficients
$\bbeta=(\beta_0,\ldots,\beta_Q)$. Different finite dictionaries can be used
when response priors are available. In this paper, we use the simplest power dictionary
$\psi_q(s)=s^q$, which gives the polynomial-linked model and makes the
rank-inflation mechanism explicit. The model is therefore not an unrestricted
nonlinear decoder and not a post-hoc calibration of a fixed Tucker fit. The
underlying signal and the response are estimated together, so the Tucker core
and factors are shaped by the nonlinear measurement model during recovery.

The key structural point is most explicit for the polynomial dictionary, where
the response creates rank inflation tied to the same Tucker factors. The degree-$p$ term
$\cS^{\od p}$ can be viewed as a Tucker tensor whose factors are polynomial
feature lifts of the original factors. Thus higher-degree channels can increase
the effective multilinear rank of the measured tensor, but they remain tied to
one low-rank Tucker signal. More general dictionaries keep the same modeling
principle: the additional expressivity must be explainable by one shared
entrywise response acting on the underlying signal. This constraint is what makes
the model useful in the intended regime and limited outside it. If the
remaining residual is ordinary multilinear structure, increasing Tucker rank is
the appropriate remedy; if the residual is coherent with a measurement
response, the learned response can provide a compact correction.

We develop this idea for masked tensor signal recovery, covering both full
reconstruction and tensor completion. The estimation procedure alternates
between a small ridge-regression update for the response coefficients and
gradient-based updates of the Tucker core and factors. To improve numerical stability,
the implementation uses normalized features and continuation when the selected
dictionary is nested.
We also introduce a response diagnostic that measures how much residual
energy can be removed by refitting only the response after a linear Tucker fit.
This diagnostic is designed to answer a practical question before full
nonlinear training: does this signal contain residual structure that is likely
to benefit from a shared entrywise response?

The resulting study is organized around the full signal-recovery pipeline:
modeling, estimation, diagnosis, and validation. On the rank side, we
characterize when an entrywise response creates efficient rank inflation and
when increasing ordinary Tucker rank is the more direct explanation. On the
estimation side, we use the special structure of the response update to keep the nonlinear
part low-dimensional. On the experimental side, we evaluate both high- and
low-diagnostic regimes to clarify the operating conditions of the model.

\subsection{Contributions}
\label{subsec:contributions}

The main contributions are summarized as follows.
\begin{itemize}
  \item We formulate nonlinear low-rank tensor recovery under an unknown
  entrywise measurement response, separating the rank of the underlying signal from
  the rank inflation observed after measurement.

  \item We propose a dictionary-linked Tucker model that keeps the Tucker ranks
  fixed while learning an entrywise response with a small number of
  coefficients. The estimator is instantiated with a power dictionary, yielding
  a compact polynomial response tied to the same Tucker core and factors.

  \item We analyze the resulting model class and estimation procedure,
  including induced rank inflation, parameter economy relative to explicit
  higher-rank Tucker models, response approximation, descent properties, diagnostic
  interpretation, and computational scaling.

  \item We evaluate the model on controlled nonlinear tensor signals and
  hyperspectral recovery tasks, with rank-inflation comparisons,
  noisy-observation robustness tests, and diagnostic stratification that
  identifies when the nonlinear response is useful.
\end{itemize}

The remainder of the paper develops the model, analysis, estimation algorithm,
diagnostic score, and experimental validation in this order.
